\documentclass[a4paper, 11pt]{book}
% ================================================================
% region Pacchetti base
% ================================================================
\usepackage[utf8]{inputenc}      % codifica UTF-8
\usepackage[T1]{fontenc}         % font T1
\usepackage[italian]{babel}      % lingua italiana

% ================================================================
% region Matematica
% ================================================================
\usepackage{amsmath, amsfonts, amssymb}
\usepackage{nicematrix}

% ================================================================
% region Grafica e diagrammi
% ================================================================
\usepackage{xcolor}
\usepackage{tikz}
\usetikzlibrary{calc, matrix}
\usepackage{pgfplots}
\pgfplotsset{compat=1.18}
\usepackage{circuitikz}
\usepackage{graphicx}

% ================================================================
% region Tabelle
% ================================================================
\usepackage{array}
\usepackage{booktabs}

% ================================================================
% region Codice sorgente
% ================================================================
\usepackage{listings}
\usepackage{fancyvrb}

\lstset{
    basicstyle=\ttfamily\small,
    keywordstyle=\color{blue!80!black}\bfseries,
    commentstyle=\color{green!50!black}\itshape,
    stringstyle=\color{orange!80!black},
    numbers=left,
    numberstyle=\tiny\color{gray},
    breaklines=true,
    showstringspaces=false,
    literate=
        {à}{{\`a}}1 {è}{{\`e}}1 {é}{{\'e}}1
        {ì}{{\`i}}1 {ò}{{\`o}}1 {ù}{{\`u}}1
        {À}{{\`A}}1 {È}{{\`E}}1 {Ì}{{\`I}}1
        {Ò}{{\`O}}1 {Ù}{{\`U}}1
}

% ================================================================
% region Box colorati
% ================================================================
\usepackage[most]{tcolorbox}
\tcbuselibrary{listings, skins, breakable}

% ================================================================
% region Layout pagina
% ================================================================
\usepackage{geometry}
\geometry{margin=2.5cm}
\setlength{\parindent}{0cm}
\setlength{\headheight}{14.49998pt}

% ================================================================
% region Varie
% ================================================================
\usepackage{enumitem}
\usepackage{comment}
\usepackage{siunitx}
\usepackage{hyperref}

% ================================================================
% region Header e piè di pagina
% ================================================================
\usepackage{fancyhdr}
\pagestyle{fancy}
\fancyhf{}
\fancyhead[L]{\nouppercase{\leftmark}}
\fancyhead[R]{\textit{prof. Bernardis Pierluigi} \qquad \qquad \thepage}
\renewcommand{\headrulewidth}{0.4pt}

\fancypagestyle{plain}{
    \fancyhf{}
    \fancyhead[L]{\nouppercase{\leftmark}}
    \fancyhead[R]{\thepage}
    \renewcommand{\headrulewidth}{0.4pt}
}

% ================================================================
% region Info autore
% ================================================================
\author{Prof. Bernardis Pierluigi}
\date{}

% ================================================================
% region Box per equazioni
% ================================================================
\tcbset{
    myequationstyle/.style={
        colback=blue!3,
        colframe=blue!100,
        arc=4pt,
        boxrule=0.8pt,
        left=6pt, right=6pt, top=4pt, bottom=4pt,
        boxsep=0pt,
        enhanced
    }
}

\newtcolorbox{eqbox}[1][]{
    myequationstyle,
    #1
}

\let\oldequation\equation
\let\endoldequation\endequation

\renewenvironment{equation}{
    \begin{eqbox}
    \begin{oldequation}
}{
    \end{oldequation}
    \end{eqbox}
}

\title{\Huge{Arduino - Componenti di Base}}

\begin{document}

\maketitle
\tableofcontents


% ================================================================
%  CAPITOLO 1 — LED
% ================================================================
\chapter{LED}

% ----------------------------------------------------------------
\section{Introduzione}
% ----------------------------------------------------------------

Il \textbf{LED} (\textit{Light Emitting Diode}, ovvero Diodo ad Emissione di Luce) è 
uno dei componenti elettronici più utilizzati al mondo. Lo troviamo ovunque nella vita 
quotidiana: nei semafori stradali, nelle lampade a risparmio energetico, negli schermi 
degli smartphone, nei display dei televisori, nelle luci dei cruscotti delle automobili 
e persino nelle fibre ottiche che trasmettono dati su internet.

Il motivo del suo enorme successo è semplice: rispetto alle tradizionali lampadine a 
incandescenza, il LED consuma molto meno energia, dura molto più a lungo (fino a 
50.000 ore di funzionamento) e si accende istantaneamente senza surriscaldarsi.

In questo capitolo impareremo cos'è un LED dal punto di vista fisico ed elettrico, 
come calcolarne correttamente il circuito di collegamento e come controllarlo con Arduino 
sia in modo digitale (acceso/spento) che analogico (variando la luminosità).

% ----------------------------------------------------------------
\section{Struttura fisica e principio di funzionamento}
% ----------------------------------------------------------------

\subsection{Cos'è un diodo}

Per capire un LED, dobbiamo prima capire cos'è un \textbf{diodo}. Un diodo è un 
componente elettronico che permette alla corrente elettrica di scorrere \textbf{in 
una sola direzione}, bloccandola nell'altra. Si comporta come una valvola a senso 
unico per la corrente.

Questa proprietà è dovuta alla struttura interna del diodo, realizzata con materiali 
\textbf{semiconduttori} (tipicamente il silicio o il germanio), ovvero materiali che 
si trovano a metà strada tra i conduttori (come il rame) e gli isolanti (come il vetro).

\subsection{Perché un LED emette luce}

Un LED è un tipo particolare di diodo che, quando la corrente lo attraversa nel verso 
corretto, emette luce. Senza entrare nei dettagli della fisica quantistica, possiamo 
capirne il funzionamento con un'analogia intuitiva.

Immaginiamo gli elettroni come palline che si muovono all'interno del materiale 
semiconduttore. Quando una pallina \textit{cade} da un livello energetico alto a uno 
basso all'interno del materiale, rilascia energia. Nei normali resistori questa energia 
si trasforma in calore. In un LED, grazie alla scelta particolare dei materiali 
semiconduttori utilizzati (arseniuro di gallio, nitruro di gallio, ecc.), questa 
energia viene rilasciata sotto forma di \textbf{fotoni}, ovvero particelle di luce.

Il colore della luce emessa dipende direttamente dal materiale semiconduttore scelto: 
materiali diversi corrispondono a livelli energetici diversi, e quindi a colori diversi.

\begin{center}
\begin{tabular}{>{\bfseries}l l l}
    \toprule
    Colore & Materiale & Tensione tipica $V_f$ \\
    \midrule
    Rosso    & Arseniuro di gallio-alluminio (AlGaAs) & \SIrange{1.8}{2.1}{\volt} \\
    Arancione & Arseniuro di gallio-fosfuro (GaAsP)   & \SIrange{2.0}{2.2}{\volt} \\
    Giallo   & Fosfuro di gallio (GaP)                & \SIrange{2.1}{2.4}{\volt} \\
    Verde    & Nitruro di gallio-indio (InGaN)         & \SIrange{2.2}{3.5}{\volt} \\
    Blu      & Nitruro di gallio (GaN)                & \SIrange{3.0}{3.5}{\volt} \\
    Bianco   & LED blu + fosforo giallo               & \SIrange{3.0}{3.5}{\volt} \\
    \bottomrule
\end{tabular}
\end{center}

\begin{tcolorbox}[colback=blue!4, colframe=blue!60, title=\textbf{Curiosità: il LED bianco}]
    La luce bianca non esiste come singolo colore: è la sovrapposizione di tutti i 
    colori dello spettro. I LED bianchi vengono costruiti combinando un LED blu ad 
    alta efficienza con uno strato di materiale fosforescente giallo. La luce blu 
    eccita il fosforo che emette luce gialla; la combinazione di blu e giallo produce 
    la percezione del bianco da parte dell'occhio umano.
\end{tcolorbox}

\subsection{Polarità del LED: anodo e catodo}

Come tutti i diodi, il LED è un componente \textbf{polarizzato}: ha un terminale 
positivo e uno negativo, e deve essere collegato nel verso giusto per funzionare.

\begin{itemize}
    \item \textbf{Anodo} (A): il terminale positivo. Si riconosce perché è la 
          \textbf{gamba più lunga} del LED.
    \item \textbf{Catodo} (K): il terminale negativo. Si riconosce perché è la 
          \textbf{gamba più corta} e perché il corpo del LED presenta una piccola 
          \textbf{piattina} sul lato del catodo.
\end{itemize}

\begin{center}
\begin{tikzpicture}[scale=1.2]
    % corpo LED
    \draw[thick, fill=gray!15] (1.5,0) -- (1.5,1.2) -- (3,0.6) -- cycle;
    \draw[thick] (3,0.2) -- (3,1.0);
    % gamba anodo (lunga)
    \draw[thick] (1.5,0.6) -- (0,0.6);
    \node[above] at (0.7,0.6) {\small Anodo (+)};
    \node[below] at (0.7,0.6) {\small gamba lunga};
    % gamba catodo (corta)
    \draw[thick] (3,0.6) -- (4.2,0.6);
    \node[above] at (3.7,0.6) {\small Catodo (-)};
    \node[below] at (3.7,0.6) {\small gamba corta};
    % piattina catodo
    \draw[thick] (3.0,0.15) -- (3.0,-0.1);
    \node[below] at (3.0,-0.1) {\tiny piattina};
    % frecce luce
    \draw[->, thick, orange!80] (2.5,1.0) -- (2.9,1.5);
    \draw[->, thick, orange!80] (2.1,1.1) -- (2.3,1.6);
    \draw[->, thick, orange!80] (1.8,0.9) -- (1.8,1.5);
\end{tikzpicture}
\end{center}

\begin{tcolorbox}[colback=red!4, colframe=red!60, title=\textbf{Attenzione — Polarità invertita}]
    Se colleghiamo il LED con la polarità invertita (anodo al negativo e catodo al 
    positivo), il LED \textbf{non si accende} e blocca il passaggio della corrente. 
    In questa condizione il LED non si danneggia, ma semplicemente non funziona. 
    Tuttavia, se la tensione inversa applicata supera il valore massimo specificato 
    nel datasheet (tipicamente \SI{5}{\volt} per LED standard), il componente 
    si danneggia irreversibilmente.
\end{tcolorbox}

% ----------------------------------------------------------------
\section{Parametri elettrici fondamentali}
% ----------------------------------------------------------------

Prima di collegare un LED a qualsiasi circuito, è indispensabile conoscere i suoi 
parametri elettrici. Questi valori si trovano nel \textbf{datasheet} del componente, 
ovvero la scheda tecnica fornita dal produttore.

\subsection{Tensione di polarizzazione diretta \texorpdfstring{$V_f$}{Vf}}

La \textbf{tensione di polarizzazione diretta} (\textit{forward voltage}, $V_f$) è 
la caduta di tensione che si misura ai capi del LED quando è acceso e la corrente 
scorre nel verso corretto. Come abbiamo visto nella tabella precedente, questo valore 
dipende dal colore del LED e si aggira tipicamente tra \SI{1.8}{\volt} e \SI{3.5}{\volt}.

È fondamentale ricordare che $V_f$ è una \textbf{proprietà fisica del componente}: 
non possiamo cambiarla. Qualunque sia la tensione di alimentazione del nostro circuito, 
il LED \textit{prenderà sempre e solo} la sua $V_f$ e nulla di più.

\subsection{Corrente diretta nominale \texorpdfstring{$I_f$}{If}}

La \textbf{corrente diretta nominale} (\textit{forward current}, $I_f$) è il valore 
di corrente per cui il LED è stato progettato per funzionare in modo ottimale e sicuro. 
Per i LED standard da 5 mm (quelli che useremo con Arduino), questo valore è 
tipicamente \SI{20}{\milli\ampere}.

\begin{tcolorbox}[colback=red!4, colframe=red!60, title=\textbf{Perché la corrente è così critica?}]
    A differenza di una resistenza, il LED \textbf{non limita autonomamente la corrente 
    che lo attraversa}. Se lo colleghiamo direttamente a una sorgente di tensione, la 
    corrente tenderebbe ad aumentare senza controllo fino a distruggere il componente 
    in pochi millisecondi. Per questo motivo è \textbf{sempre obbligatorio} inserire 
    una resistenza di limitazione in serie al LED.
\end{tcolorbox}

\subsection{Corrente massima assoluta}

La \textbf{corrente massima assoluta} è il valore oltre il quale il LED si danneggia 
in modo permanente e istantaneo. Per i LED standard è circa \SI{30}{\milli\ampere}. 
Non si deve mai avvicinarsi a questo valore: in pratica, per un uso sicuro con Arduino, 
è consigliabile lavorare con correnti tra \SI{5}{\milli\ampere} e \SI{15}{\milli\ampere}.

\begin{tcolorbox}[colback=blue!4, colframe=blue!60, title=\textbf{Nota pratica}]
    I pin digitali di Arduino Uno possono erogare al massimo \SI{40}{\milli\ampere}, 
    ma la scheda nel suo complesso non dovrebbe superare \SI{200}{\milli\ampere} totali. 
    Per sicurezza, è buona norma progettare i circuiti per non superare i 
    \SI{10}{\milli\ampere} per LED.
\end{tcolorbox}

% ----------------------------------------------------------------
\section{La resistenza di limitazione}
% ----------------------------------------------------------------

\subsection{Perché è necessaria}

Abbiamo visto che il LED non può essere collegato direttamente alla tensione di 
alimentazione perché la corrente diventerebbe eccessiva. La soluzione è inserire 
in \textbf{serie} al LED una \textbf{resistenza di limitazione} $R$, il cui compito 
è quello di \textit{assorbire} la differenza di tensione tra l'alimentazione e la 
$V_f$ del LED, limitando di conseguenza la corrente.

\subsection{Calcolo con la legge di Kirchhoff e la legge di Ohm}

Consideriamo il circuito in serie composto da: sorgente di tensione $V_{cc}$, 
resistenza $R$ e LED con tensione $V_f$.

Applicando la \textbf{legge di Kirchhoff alle tensioni} (la somma delle tensioni in 
un anello chiuso è zero), otteniamo:

\begin{equation}
    V_{cc} = V_R + V_f
\end{equation}

Dove $V_R$ è la caduta di tensione sulla resistenza. Ricaviamo $V_R$:

\begin{equation}
    V_R = V_{cc} - V_f
\end{equation}

In un circuito in serie la corrente è uguale in ogni punto, quindi la corrente che 
attraversa la resistenza è la stessa che attraversa il LED: $I_R = I_f$. 
Applicando la \textbf{legge di Ohm} alla resistenza ($V = R \cdot I$):

\begin{equation}
    R = \frac{V_R}{I_f} = \frac{V_{cc} - V_f}{I_f}
\end{equation}

\subsection{Esempio di calcolo}

Vogliamo collegare un \textbf{LED rosso} al pin 13 di Arduino, che fornisce 
\SI{5}{\volt}. Il LED rosso ha $V_f = \SI{2.0}{\volt}$ e vogliamo farlo 
lavorare a $I_f = \SI{10}{\milli\ampere}$ per sicurezza.

\begin{equation}
    R = \frac{V_{cc} - V_f}{I_f} = \frac{\SI{5}{\volt} - \SI{2.0}{\volt}}{\SI{10}{\milli\ampere}} = \frac{\SI{3}{\volt}}{\SI{0.010}{\ampere}} = \SI{300}{\ohm}
\end{equation}

Il valore esatto di \SI{300}{\ohm} potrebbe non essere disponibile nella serie 
standard di resistenze commerciali. In questo caso scegliamo il valore 
\textbf{immediatamente superiore} disponibile, cioè \textbf{\SI{330}{\ohm}}: 
una resistenza leggermente più alta riduce leggermente la corrente (e quindi la 
luminosità), ma protegge meglio il componente. Non si sceglie mai un valore 
inferiore, perché aumenterebbe la corrente oltre il valore di progetto.

\begin{tcolorbox}[colback=green!4, colframe=green!60, title=\textbf{Verifica del calcolo}]
    Con $R = \SI{330}{\ohm}$, la corrente effettiva sarà:
    \[
        I_f = \frac{V_{cc} - V_f}{R} = \frac{\SI{5}{\volt} - \SI{2.0}{\volt}}{\SI{330}{\ohm}} \approx \SI{9.1}{\milli\ampere}
    \]
    Valore ben al di sotto dei \SI{20}{\milli\ampere} nominali: il LED è al sicuro.
\end{tcolorbox}

\subsection{Leggere il valore di una resistenza: il codice a colori}

Le resistenze fisiche non hanno il valore scritto sopra (sarebbero troppo piccole), 
ma utilizzano un \textbf{codice a colori} di bande colorate. Per una resistenza a 
4 bande:

\begin{center}
\begin{tabular}{>{\bfseries}l c c}
    \toprule
    Colore   & Cifra & Moltiplicatore \\
    \midrule
    Nero     & 0 & $\times 1$ \\
    Marrone  & 1 & $\times 10$ \\
    Rosso    & 2 & $\times 100$ \\
    Arancione & 3 & $\times 1\,000$ \\
    Giallo   & 4 & $\times 10\,000$ \\
    Verde    & 5 & $\times 100\,000$ \\
    Blu      & 6 & $\times 1\,000\,000$ \\
    Viola    & 7 & --- \\
    Grigio   & 8 & --- \\
    Bianco   & 9 & --- \\
    Oro      & --- & $\times 0.1$ (tolleranza $\pm5\%$) \\
    Argento  & --- & $\times 0.01$ (tolleranza $\pm10\%$) \\
    \bottomrule
\end{tabular}
\end{center}

Una resistenza da \SI{330}{\ohm} ha le bande: \textbf{arancione -- arancione -- marrone -- oro} 
(3 -- 3 -- $\times$10 -- $\pm$5\% = 330\,\si{\ohm}).

% ----------------------------------------------------------------
\section{Schema di collegamento}
% ----------------------------------------------------------------

\subsection{Componenti necessari}

\begin{itemize}
    \item 1$\times$ Arduino Uno (o compatibile)
    \item 1$\times$ LED (colore a scelta)
    \item 1$\times$ Resistenza da \SI{330}{\ohm}
    \item 1$\times$ Breadboard
    \item 2 cavi jumper maschio-maschio
\end{itemize}

\subsection{Schema elettrico}

\begin{center}
\begin{circuitikz}[scale=1.1, american]
    \draw (0,0) node[left] {\texttt{Pin 13}}
        to[short] (1,0)
        to[R, l=\SI{330}{\ohm}] (3.5,0)
        to[led, l=LED] (5.5,0)
        to[short] (6.5,0) node[ground] {};
    \node[below] at (6.5,-0.3) {\texttt{GND}};
\end{circuitikz}
\end{center}

\subsection{Procedura di montaggio su breadboard}

La breadboard è una basetta forata che permette di realizzare circuiti elettronici 
senza saldature. I fori sono collegati internamente in file orizzontali (al centro) 
e in colonne verticali (sui bordi, riservate all'alimentazione).

\begin{enumerate}
    \item Collega un cavo jumper dal \textbf{pin 13} di Arduino a una riga 
          della breadboard (es. riga 10, colonna e).
    \item Inserisci la \textbf{resistenza da \SI{330}{\ohm}} con un terminale 
          nella stessa riga del cavo (es. riga 10) e l'altro terminale in 
          una riga diversa (es. riga 10, saltando la separazione centrale 
          se necessario).
    \item Inserisci il \textbf{LED} con l'\textbf{anodo} (gamba lunga) nella 
          stessa colonna dell'altro terminale della resistenza e il 
          \textbf{catodo} (gamba corta) in una riga adiacente libera.
    \item Collega un cavo jumper dal catodo del LED al pin \textbf{GND} di Arduino.
\end{enumerate}

\begin{tcolorbox}[colback=orange!4, colframe=orange!70, title=\textbf{Regola d'oro della breadboard}]
    I fori sulla stessa riga orizzontale (stessa lettera, stesso numero) 
    \textbf{non} sono collegati tra loro. I fori sulla stessa colonna verticale 
    (stesso numero, lettere a-e oppure f-j) \textbf{sono} collegati tra loro. 
    Controlla sempre bene prima di alimentare il circuito.
\end{tcolorbox}

% ----------------------------------------------------------------
\section{Programmazione con Arduino}
% ----------------------------------------------------------------

\subsection{Struttura di uno sketch Arduino}

Un programma per Arduino si chiama \textbf{sketch} e ha sempre la stessa struttura 
di base, con due funzioni obbligatorie:

\begin{itemize}
    \item \texttt{setup()}: viene eseguita \textbf{una sola volta} all'avvio della 
          scheda o dopo ogni reset. Qui si configurano i pin e si inizializzano 
          le variabili.
    \item \texttt{loop()}: viene eseguita \textbf{in modo continuo e ripetuto} 
          dopo la \texttt{setup()}. Contiene il comportamento principale del programma.
\end{itemize}

\subsection{Configurare un pin come uscita: \texttt{pinMode()}}

Prima di poter controllare un LED (o qualsiasi altro componente), dobbiamo dire 
ad Arduino se quel pin deve comportarsi da \textbf{uscita} (Arduino manda tensione 
al circuito) o da \textbf{ingresso} (Arduino legge la tensione dal circuito). 
Questo si fa con la funzione \texttt{pinMode()}:

\begin{center}
\texttt{pinMode(numeroPIN, modalità);}
\end{center}

\begin{itemize}
    \item \texttt{numeroPIN}: il numero del pin da configurare (es. \texttt{13})
    \item \texttt{modalità}: \texttt{OUTPUT} per uscita, \texttt{INPUT} per ingresso
\end{itemize}

\subsection{Accendere e spegnere il LED: \texttt{digitalWrite()}}

Per portare un pin digitale a \SI{5}{\volt} (LED acceso) o a \SI{0}{\volt} 
(LED spento) si usa \texttt{digitalWrite()}:

\begin{center}
\texttt{digitalWrite(numeroPIN, valore);}
\end{center}

\begin{itemize}
    \item \texttt{HIGH}: porta il pin a \SI{5}{\volt} → LED \textbf{acceso}
    \item \texttt{LOW}: porta il pin a \SI{0}{\volt} → LED \textbf{spento}
\end{itemize}

\subsection{Mettere in pausa il programma: \texttt{delay()}}

La funzione \texttt{delay(ms)} mette in pausa l'esecuzione del programma per il 
numero di millisecondi specificato. Per esempio, \texttt{delay(1000)} aspetta 
esattamente 1 secondo prima di proseguire.

\subsection{Esempio 1 — Blink (lampeggio)}

Il programma \textit{Blink} è il classico ``Hello World'' di Arduino: fa lampeggiare 
il LED con un ritmo regolare di 1 secondo acceso e 1 secondo spento.

\begin{tcblisting}{
    listing only,
    listing engine=listings,
    colback=gray!5,
    colframe=gray!50,
    title=\textbf{Esempio 1 — Blink},
    breakable,
    listing options={
        language=C++,
        basicstyle=\ttfamily\small,
        keywordstyle=\color{blue!80!black}\bfseries,
        commentstyle=\color{green!50!black}\itshape,
        stringstyle=\color{orange!80!black},
        numbers=left,
        numberstyle=\tiny\color{gray},
        breaklines=true,
        showstringspaces=false,
        morekeywords={pinMode, digitalWrite, delay, OUTPUT, HIGH, LOW, void, const, int}
    }
}
// Definiamo una costante per il numero del pin.
// Usando una costante (invece di scrivere 13 direttamente),
// se in futuro cambiamo pin, basta modificare solo questa riga.
const int PIN_LED = 13;

void setup() {
  // Configuriamo PIN_LED come uscita digitale.
  // Questo dice ad Arduino: "da questo pin voglio inviare tensione".
  pinMode(PIN_LED, OUTPUT);
}

void loop() {
  digitalWrite(PIN_LED, HIGH);  // Accende il LED (porta il pin a 5V)
  delay(1000);                  // Aspetta 1 secondo (1000 millisecondi)
  digitalWrite(PIN_LED, LOW);   // Spegne il LED (porta il pin a 0V)
  delay(1000);                  // Aspetta 1 secondo
  // La funzione loop() ricomincia automaticamente dall'inizio:
  // il LED continuerà a lampeggiare indefinitamente.
}
\end{tcblisting}

\textbf{Analisi del codice riga per riga:}
\begin{itemize}
    \item \textbf{Riga 6} — \texttt{const int PIN\_LED = 13;}: dichiariamo una 
          costante intera di nome \texttt{PIN\_LED} con valore 13. La parola chiave 
          \texttt{const} significa che il valore non potrà essere modificato 
          durante l'esecuzione del programma.
    \item \textbf{Riga 11} — \texttt{pinMode(PIN\_LED, OUTPUT);}: eseguita una 
          sola volta all'avvio, configura il pin 13 come uscita digitale.
    \item \textbf{Riga 16} — \texttt{digitalWrite(PIN\_LED, HIGH);}: porta il 
          pin 13 a \SI{5}{\volt}. La corrente scorre attraverso la resistenza 
          e il LED, che si accende.
    \item \textbf{Riga 17} — \texttt{delay(1000);}: il programma si blocca per 
          1000 ms. Durante questo tempo il LED rimane nello stato in cui si trova.
    \item \textbf{Riga 18} — \texttt{digitalWrite(PIN\_LED, LOW);}: porta il 
          pin 13 a \SI{0}{\volt}. Nessuna differenza di tensione, nessuna corrente, 
          LED spento.
\end{itemize}

\subsection{Variare la luminosità: PWM e \texttt{analogWrite()}}

I pin digitali di Arduino possono essere solo in due stati: \SI{5}{\volt} (HIGH) 
oppure \SI{0}{\volt} (LOW). Non è possibile impostare direttamente una tensione 
intermedia di, ad esempio, \SI{2.5}{\volt}. Come facciamo allora a variare la 
luminosità del LED?

La soluzione si chiama \textbf{PWM} (\textit{Pulse Width Modulation}, Modulazione 
di Larghezza di Impulso). L'idea è la seguente: invece di tenere il pin fisso a 
un valore di tensione intermedio, il pin viene commutato rapidamente tra HIGH e LOW 
a una frequenza di circa \SI{490}{\hertz} (490 volte al secondo). L'occhio umano, 
non riuscendo a percepire variazioni così rapide, vede una luminosità \textit{media} 
proporzionale al tempo in cui il pin è rimasto HIGH.

Questo rapporto tra il tempo in cui il segnale è HIGH e il periodo totale si chiama 
\textbf{duty cycle}:

\begin{equation}
    \text{Duty Cycle} = \frac{t_{HIGH}}{T} \times 100\%
\end{equation}

Un duty cycle del 50\% significa che il pin è HIGH per metà del tempo: il LED 
appare a metà della sua luminosità massima. Un duty cycle del 100\% equivale a 
\texttt{digitalWrite(pin, HIGH)}: LED alla massima luminosità.

\begin{center}
\begin{tikzpicture}[scale=0.85]
    % Duty cycle 25%
    \node[left] at (0,3.4) {\small 25\%};
    \draw[thick, blue] (0,3) -- (0.5,3) -- (0.5,3.8) -- (1,3.8) -- (1,3) -- (2,3) -- (2,3.8) -- (2.5,3.8) -- (2.5,3) -- (4,3);
    % Duty cycle 50%
    \node[left] at (0,1.9) {\small 50\%};
    \draw[thick, blue] (0,1.5) -- (1,1.5) -- (1,2.3) -- (2,2.3) -- (2,1.5) -- (3,1.5) -- (3,2.3) -- (4,2.3) -- (4,1.5);
    % Duty cycle 75%
    \node[left] at (0,0.4) {\small 75\%};
    \draw[thick, blue] (0,0) -- (1.5,0) -- (1.5,0.8) -- (2,0.8) -- (2,0) -- (3.5,0) -- (3.5,0.8) -- (4,0.8);
    % titolo
    \node[right] at (4.2,3.4) {\small luminosità bassa};
    \node[right] at (4.2,1.9) {\small luminosità media};
    \node[right] at (4.2,0.4) {\small luminosità alta};
\end{tikzpicture}
\end{center}

La funzione \texttt{analogWrite(pin, valore)} imposta il duty cycle del pin PWM 
specificato. Il parametro \texttt{valore} va da \textbf{0} (duty cycle 0\%, LED 
completamente spento) a \textbf{255} (duty cycle 100\%, LED alla massima 
luminosità).

\begin{tcolorbox}[colback=orange!4, colframe=orange!70, title=\textbf{Quali pin supportano il PWM?}]
    Su Arduino Uno, \textbf{non tutti i pin} supportano il PWM. I pin abilitati 
    al PWM sono contrassegnati dal simbolo \textbf{\texttildelow} sulla scheda e 
    sono: \textbf{3, 5, 6, 9, 10, 11}. Se si tenta di usare \texttt{analogWrite()} 
    su un pin non PWM, il comportamento è imprevedibile. Per questa ragione, 
    negli esempi che usano il PWM, sposteremo il LED al pin 9.
\end{tcolorbox}

\subsection{Esempio 2 — Fade (dissolvenza)}

\begin{tcblisting}{
    listing only,
    listing engine=listings,
    colback=gray!5,
    colframe=gray!50,
    title=\textbf{Esempio 2 — Fade},
    breakable,
    listing options={
        language=C++,
        basicstyle=\ttfamily\small,
        keywordstyle=\color{blue!80!black}\bfseries,
        commentstyle=\color{green!50!black}\itshape,
        stringstyle=\color{orange!80!black},
        numbers=left,
        numberstyle=\tiny\color{gray},
        breaklines=true,
        showstringspaces=false,
        morekeywords={pinMode, analogWrite, delay, OUTPUT, void, const, int, for}
    }
}
// Usiamo il pin 9 perché supporta il PWM (simbolo ~ sulla scheda)
const int PIN_LED = 9;

void setup() {
  pinMode(PIN_LED, OUTPUT);
}

void loop() {
  // --- Fase 1: aumenta gradualmente la luminosità ---
  // La variabile 'luce' parte da 0 e arriva a 255, incrementandosi di 1.
  // Ad ogni passo impostiamo il duty cycle del pin e aspettiamo 10 ms.
  for (int luce = 0; luce <= 255; luce++) {
    analogWrite(PIN_LED, luce);
    delay(10);   // 10 ms × 256 passi = ~2.5 secondi per salire
  }

  // --- Fase 2: diminuisce gradualmente la luminosità ---
  for (int luce = 255; luce >= 0; luce--) {
    analogWrite(PIN_LED, luce);
    delay(10);   // altri ~2.5 secondi per scendere
  }
  // Il loop riparte: effetto di respiro continuo.
}
\end{tcblisting}

\textbf{Analisi del ciclo \texttt{for}:}

Il ciclo \texttt{for} ha la struttura:
\begin{center}
    \texttt{for (inizializzazione; condizione; incremento) \{ ... \}}
\end{center}
\begin{itemize}
    \item \textbf{Inizializzazione} (\texttt{int luce = 0}): crea la variabile 
          \texttt{luce} e la imposta a 0. Viene eseguita \textit{una sola volta} 
          all'inizio del ciclo.
    \item \textbf{Condizione} (\texttt{luce <= 255}): il ciclo continua finché 
          questa condizione è vera. Quando \texttt{luce} supera 255, il ciclo termina.
    \item \textbf{Incremento} (\texttt{luce++}): equivale a \texttt{luce = luce + 1}. 
          Viene eseguito alla fine di ogni iterazione.
\end{itemize}

\subsection{Esempio 3 — Lampeggio a frequenza variabile con potenziometro}

Questo esempio anticipa il prossimo capitolo sul potenziometro, ma mostra come 
rendere il lampeggio interattivo. La velocità di lampeggio viene controllata 
da una lettura analogica:

\begin{tcblisting}{
    listing only,
    listing engine=listings,
    colback=gray!5,
    colframe=gray!50,
    title=\textbf{Esempio 3 — Blink a velocità variabile},
    breakable,
    listing options={
        language=C++,
        basicstyle=\ttfamily\small,
        keywordstyle=\color{blue!80!black}\bfseries,
        commentstyle=\color{green!50!black}\itshape,
        stringstyle=\color{orange!80!black},
        numbers=left,
        numberstyle=\tiny\color{gray},
        breaklines=true,
        showstringspaces=false,
        morekeywords={pinMode, digitalWrite, analogRead, map, delay, OUTPUT, HIGH, LOW, void, const, int, long}
    }
}
const int PIN_LED  = 13;
const int PIN_POT  = A0;  // Pin analogico per il potenziometro

void setup() {
  pinMode(PIN_LED, OUTPUT);
}

void loop() {
  // Leggiamo il valore del potenziometro: restituisce un intero tra 0 e 1023
  int valPot = analogRead(PIN_POT);

  // Convertiamo il valore in un intervallo utile: da 50 ms a 1000 ms
  long intervallo = map(valPot, 0, 1023, 50, 1000);

  // Lampeggio con il tempo determinato dal potenziometro
  digitalWrite(PIN_LED, HIGH);
  delay(intervallo);
  digitalWrite(PIN_LED, LOW);
  delay(intervallo);
}
\end{tcblisting}

La funzione \texttt{map(valore, minIn, maxIn, minOut, maxOut)} converte linearmente 
un valore da un intervallo di ingresso a un intervallo di uscita. In questo caso, 
il valore grezzo del potenziometro (0--1023) viene convertito in un tempo di 
attesa (50--1000 ms).

% ----------------------------------------------------------------
\section{Applicazioni pratiche}
% ----------------------------------------------------------------

I LED non servono solo come segnalatori luminosi elementari. Ecco alcune applicazioni 
concrete che potrete realizzare combinando quanto appreso in questo capitolo:

\begin{itemize}
    \item \textbf{Spia di stato}: un LED verde fisso indica che il sistema è 
          operativo; un LED rosso lampeggiante segnala un errore o un allarme.
    \item \textbf{Indicatore di livello}: una serie di LED disposti in fila, 
          accesi in sequenza, simulano una barra di avanzamento (come la 
          batteria di uno smartphone).
    \item \textbf{Effetto \textit{breathing}}: variando il duty cycle con una 
          funzione sinusoidale anziché lineare, si ottiene un effetto di 
          ``respiro'' molto più morbido e naturale.
    \item \textbf{Comunicazione ottica}: modulando il LED a frequenze elevate, 
          è possibile trasmettere dati binari. Questo è il principio alla base 
          del telecomando a infrarossi.
    \item \textbf{LED RGB}: i LED RGB contengono tre diodi (rosso, verde, blu) 
          in un unico package. Controllando separatamente il duty cycle di 
          ciascun colore, è possibile ottenere qualsiasi colore dello spettro 
          visibile.
\end{itemize}

% ----------------------------------------------------------------
\section{Domande di verifica}
% ----------------------------------------------------------------

\begin{enumerate}
    \item Cosa significa la sigla LED e da quali componenti fisici è costituito?
    
    \item Come si riconosce fisicamente l'anodo dal catodo di un LED? 
          Descrivi almeno due metodi.
    
    \item Un LED viene collegato direttamente ai \SI{5}{\volt} di Arduino senza 
          resistenza. Cosa succede e perché?
    
    \item Calcola la resistenza di limitazione necessaria per collegare ad Arduino 
          un LED \textbf{blu} con $V_f = \SI{3.2}{\volt}$, volendo una corrente 
          di lavoro di $\SI{8}{\milli\ampere}$. Quale valore commerciale sceglieresti?
    
    \item Quali bande colorate ha una resistenza da \SI{470}{\ohm} con tolleranza 
          del 5\%?
    
    \item Qual è la differenza tra la funzione \texttt{digitalWrite()} e la 
          funzione \texttt{analogWrite()}? In quale situazione useresti l'una 
          e in quale l'altra?
    
    \item Cosa si intende per \textbf{duty cycle}? Se impostassimo 
          \texttt{analogWrite(9, 128)}, a quale percentuale di duty cycle 
          corrisponde (approssimativamente)?
    
    \item Un LED è collegato al pin 6 di Arduino (PWM). Lo sketch contiene 
          \texttt{analogWrite(6, 64)}. Stima la luminosità percepita rispetto 
          alla massima e spiega il ragionamento.
    
    \item Nel codice dell'Esempio 1, cosa succederebbe se rimuovessimo la seconda 
          \texttt{delay(1000)} (quella dopo \texttt{LOW})? Descrivi il 
          comportamento del LED.
    
    \item Modifica l'Esempio 2 in modo che il LED raggiunga la massima luminosità 
          in 1 secondo (anziché 2.5 secondi). Mostra la modifica al codice e 
          spiega il ragionamento.
\end{enumerate}








\chapter{Pulsante}
\chapter{Potenziometro}
\chapter{Fotoresistore}
\chapter{Sensore di Distanza HC-SR04}



\end{document}